% Legend Page
% ====================================================================
\newpage
\section*{Legend: Box-Type Reference}
\addcontentsline{toc}{section}{Legend: Box-Type Reference}

All mathematical content in these notes is displayed in colour-coded boxes.
The table below gives the full reference; subsequent new content must follow
this convention.

\begin{center}
\renewcommand{\arraystretch}{1.35}
\begin{tabular}{>{\ttfamily}l l p{7.5cm}}
\toprule
\textrm{\textbf{Environment}} & \textbf{Colour} & \textbf{Usage} \\
\midrule
defbox    & \textcolor{defblue}{Blue}                 & Definitions; formal notation \\
thmbox    & \textcolor{thmcrimson}{Crimson}           & Theorems; Lemmas \\
proofbox  & \textcolor{proofgray}{Gray}               & Full or sketched proofs \\
corbox    & \textcolor{corteal}{Teal}                 & Corollaries \\
propbox   & \textcolor{propviolet}{Violet}            & Propositions \\
tipbox    & \textcolor{tipamber}{Amber}               & Remarks; Notes; Recalls \\
warnbox   & \textcolor{warnred}{Red}                  & Warnings; important notices; exam scope \\
flowbox   & \textcolor{orange!65!black}{Dark orange}  & Key formulae; summary boxes; recurrences \\
exbox     & \textcolor{orange!40!brown}{Warm umber}   & Worked examples \\
methodbox & \textcolor{methodsteel}{Steel blue}       & Proof strategies; algorithms; methods \\
\bottomrule
\end{tabular}
\end{center}

\begin{tipbox}[Remark --- Primary Reference]
K.\ Hoffman and R.\ Kunze, \textit{Linear Algebra}, 2nd ed.\ (Prentice Hall).
These notes follow the lecture structure and are not tied to any particular
chapter ordering of the textbook.
\end{tipbox}